\section{Introduction}

\subsection{Source Code}

The source code for this project is available in a public GitHub repository at
\url{https://github.com/Hcapacity/IOT_Project.git}. The repository contains the complete implementation of the ESP32-S3 application, including all source files, configuration files, and documentation.
\subsection{Project Context}
The objective of this project is to transform a baseline RTOS-based embedded
application into a complete Internet of Things system on the ESP32-S3 platform.
The target board provides sufficient GPIO, Wi-Fi connectivity, and compute
resources to host multiple concurrent tasks under FreeRTOS while still keeping
the implementation close to the constraints of an embedded device. The project
is built with PlatformIO and the Arduino framework, but its internal
organization follows a classical RTOS design with explicitly created tasks,
queues, semaphores, and mutexes rather than a monolithic \texttt{loop()}
application. This is important because the assignment is not only about adding
features; it also requires students to demonstrate understanding of task
coordination, modularity, and synchronization in an embedded real-time system
\cite{freertos_kernel,platformio,esp32_arduino}.

The final system can be understood as a smart environmental node. It reads
temperature and humidity, presents the state locally through LEDs and an LCD,
offers a browser-based configuration and monitoring interface, performs
on-device TinyML inference, and publishes telemetry to a cloud platform.
Compared to a simple sensor board, the project therefore integrates four
different interaction layers:
\begin{enumerate}
  \item the physical sensing layer (DHT20, LCD, LEDs, NeoPixel),
  \item the RTOS coordination layer (tasks, queues, semaphores, shared state),
  \item the local network user-interface layer (AP/STA web server), and
  \item the cloud-connected telemetry layer (CoreIOT via MQTT).
\end{enumerate}

\subsection{Assignment Tasks and Design Implications}
The assignment defines six major tasks. Understanding the constraints of each
task is essential because the final architecture must support all of them
simultaneously instead of implementing them in isolation.

\subsubsection*{Task 1: Single LED blink with temperature conditions.}
The first requirement is to redefine the blink behavior of a single LED
according to at least three temperature states. This requirement looks simple,
but it introduces two important design issues. First, the LED task must not
block sensor acquisition or networking. Second, LED mode changes should happen
immediately when a new temperature value arrives. This naturally suggests a
dedicated LED manager task that receives temperature-derived commands through a
queue.

\subsubsection*{Task 2: NeoPixel control based on humidity.}
The second task requires humidity-dependent color feedback. In this project,
the NeoPixel is not used as a static indicator but as a continuous color mapper
from low to high humidity. This expands the assignment requirement into a more
expressive visual representation. Because humidity values are produced by the
sensor task and consumed asynchronously by the NeoPixel task, queue-based data
delivery is again the most appropriate mechanism.

\subsubsection*{Task 3: Temperature and humidity monitoring with LCD display.}
The LCD requirement is more subtle than it first appears. The display must show
environmental information, react to system state, and demonstrate RTOS
synchronization. In the implemented system, the LCD is not limited to one
screen. It supports sensor display, Wi-Fi status display, and TinyML forecast
display. The task also forces correct handling of a shared I\textsuperscript{2}C
bus, which is why the code uses a mutex to protect both the DHT20 and LCD
access.

\subsubsection*{Task 4: Web server in Access Point mode.}
The web server task introduces a complete user-interface requirement on top of
the embedded node. The board must expose a meaningful AP-mode interface and
allow control of two actions. In this project, the AP dashboard not only shows
live sensor data but also supports Wi-Fi scanning, Wi-Fi profile storage, AP to
STA switching, external LED on/off control, and external LED brightness
control. This task directly influences the architecture because the web task
must remain responsive while other RTOS tasks continue to run periodically.

\subsubsection*{Task 5: TinyML deployment and accuracy evaluation.}
This task requires more than simply executing a model. The system must define a
data pipeline, prepare model inputs, deploy inference on the ESP32-S3, and
evaluate the result. In the current implementation, TinyML is integrated as a
dedicated consumer task that accumulates a batch of sensor readings, derives
eight normalized input features, invokes a TensorFlow Lite Micro model, and
publishes the result to the rest of the system. This requirement has a strong
architectural impact because inference is periodic, compute-intensive relative
to the rest of the application, and semantically different from raw sensing.

\subsubsection*{Task 6: Data publishing to CoreIOT cloud server.}
Cloud integration introduces the requirement that the board must publish data in
STA mode only, with the correct authentication token. In practice, this means
the application needs a mechanism to detect when Wi-Fi has successfully moved
from AP or disconnected state into valid STA connectivity. This project
implements that transition with a binary semaphore that signals ``Internet
ready'' from the web task to the CoreIOT task. The cloud task therefore does
not need to continuously probe for connectivity from boot time.

\subsection{From Requirements to System Construction}
Taken together, the six assignment tasks strongly suggest a distributed task
architecture rather than a sequential application. A single-threaded design
would quickly become difficult to maintain, because sensor acquisition, LCD
rendering, blink control, Wi-Fi handling, HTTP requests, TinyML inference, and
MQTT publishing all have different execution patterns and timing expectations.

For this reason, the implemented system is built around seven FreeRTOS tasks:
\texttt{sensor\_task}, \texttt{led\_manager\_task},
\texttt{neo\_pixel\_task}, \texttt{lcd\_task}, \texttt{main\_server\_task},
\texttt{coreiot\_task}, and \texttt{tiny\_ml\_task}. The main function only
allocates an application context, creates queues and synchronization primitives,
and starts these tasks. Each task is intentionally narrow in responsibility:
the sensor task owns periodic sensing, the LED and NeoPixel tasks own local
actuation, the LCD task owns presentation, the web task owns UI and Wi-Fi mode
management, the TinyML task owns prediction, and the CoreIOT task owns cloud
communication.

This decomposition also makes the system easier to explain and test. Each
assignment task can be traced to one or more RTOS tasks and their data
interfaces. For example, the sensor task produces the physical environment
measurements, the LED and NeoPixel tasks consume them to fulfill Tasks 1 and 2,
the LCD and web tasks consume them to fulfill Tasks 3 and 4, the TinyML task
transforms them to fulfill Task 5, and the CoreIOT task publishes them to
fulfill Task 6.

\subsection{Project Objectives}
Based on the assignment and the implemented architecture, the practical
objectives of this project are the following:
\begin{itemize}
  \item to build a stable multi-task IoT application on ESP32-S3;
  \item to demonstrate correct use of FreeRTOS queues, semaphores, and mutexes;
  \item to implement meaningful local feedback for temperature and humidity;
  \item to provide a usable AP-mode dashboard for monitoring and control;
  \item to support STA-mode cloud publishing to CoreIOT;
  \item to integrate TinyML inference into the live sensing pipeline; and
  \item to keep the design modular enough for future extension.
\end{itemize}

\subsection{Contributions of the Implemented System}
Compared with a minimal RTOS demonstration, the current system contributes the
following integrated features:
\begin{itemize}
  \item periodic sensing every 2 seconds using a dedicated sensor task;
  \item temperature-dependent LED blinking with multiple operating modes;
  \item humidity-dependent NeoPixel color mapping;
  \item a multi-mode LCD that adapts to AP, STA, and TinyML states;
  \item a browser-accessible AP dashboard with charting and device control;
  \item Wi-Fi scanning, profile storage, and automatic reconnection flow;
  \item batched TinyML inference using an 8-feature input vector;
  \item CoreIOT MQTT telemetry publishing every 10 seconds in STA mode; and
  \item a shared application context that reduces the need for project-wide
    global variables.
\end{itemize}

The remainder of this report explains how these features map to the assignment,
how they are implemented in code, and how the behavior of each task can be
analyzed from a system-design and RTOS perspective.
